%% Full reciprocal ZF proofs for the finite BPI/ultrafilter supplement.
\begin{annexitem}{annex:p01-bpi-duality}{thm:p01-bpi-duality}{Choice-free filter and ideal calculus}
\begin{proof}
For \(E\subseteq B\), put
\[
 F_E=\{b\in B:\text{for some }n\in\mathbb N\text{ and }a_1,\ldots,a_n\in E,
                    \ \bigwedge_{k=1}^{n}a_k\leq_B b\}.
\]
The empty meet gives \(1_B\in F_E\); a one-member meet gives
\(E\subseteq F_E\). Upward closure follows from transitivity of
\(\leq_B\). If two finite meets witness \(b,c\in F_E\), their
combined finite list witnesses \(b\wedge c\in F_E\). Hence \(F_E\)
is a filter. Every filter containing \(E\) contains each such finite
meet and everything above it, so contains \(F_E\). Finally,
\(0_B\in F_E\) iff one witnessing meet is zero, proving the
properness criterion, including the empty-list case.

For clarity, the Boolean identities used next follow from the declared
axioms, without a representation theorem. If \(r\wedge s=0_B\), then
\(r=r\wedge(s\vee\neg s)=r\wedge\neg s\), so
\(r\leq_B\neg s\); the converse follows by meeting with \(s\).
Thus \(a\leq_B b\) implies \(\neg b\leq_B\neg a\).
Complements are unique: if \(x,y\) both complement \(a\), then
\(x=x\wedge(a\vee y)=x\wedge y\), and symmetrically
\(y=x\wedge y\). Distribution verifies that \(\neg a\wedge\neg b\)
has zero meet and total join one with \(a\vee b\); it is therefore
\(\neg(a\vee b)\). Double complementation follows from uniqueness,
and gives the dual De Morgan identity. Consequently complement reverses
Boolean order and interchanges meets and joins. Thus
\(d(F)\) contains \(0_B\); if \(\neg a\in d(F)\) and
\(b\leq_B\neg a\), then \(a\leq_B\neg b\), so
\(\neg b\in F\) and \(b\in d(F)\). Also
\(\neg a\vee\neg c=\neg(a\wedge c)\in d(F)\) for
\(a,c\in F\). Therefore \(d(F)\) is an ideal. The same argument
with the operations interchanged takes ideals to filters.
Since \(d(d(E))=E\) and \(d\) preserves inclusion, this is the
stated bijection. Moreover \(1_B\in d(F)\) iff \(0_B\in F\).
Applying this bijection to the generated-filter argument proves that the
ideal generated by \(E\) is exactly
\(\{b:b\leq_B\bigvee_{k=1}^{n}a_k\text{ for a finite list in }E\}\),
with empty join \(0_B\).

Let \(U\) be a maximal proper filter. It cannot contain both
\(a\) and \(\neg a\), since their meet is zero. If \(a\notin U\),
the filter generated by \(U\cup\{a\}\) is improper by maximality.
Combining the finitely many members from \(U\) in a witnessing list
gives \(u\in U\) with \(u\wedge a=0_B\). (A witness not using
\(a\) would contradict properness.) The identity
\(u=(u\wedge a)\vee(u\wedge\neg a)\) implies
\(u\leq_B\neg a\), and hence \(\neg a\in U\).
Thus \(U\) decides complements. Conversely, if a proper filter decides
complements, a larger proper filter containing some \(a\notin U\)
would also contain \(\neg a\in U\), an impossibility. It is maximal.

Decision implies filter primeness: if \(a,b\notin U\), then
\(\neg a,\neg b\in U\), so
\(\neg(a\vee b)\in U\) and \(a\vee b\notin U\).
Conversely, filter primeness applied to
\(a\vee\neg a=1_B\in U\), together with properness, gives
exactly-one decision. For any filter,
\(a\wedge b\in U\) iff both \(a,b\in U\), by closure and
upward closure. Consequently decision and primeness give all the equations
\[
\begin{aligned}
 \chi_U(0_B)&=0,&\chi_U(1_B)&=1,\\
 \chi_U(\neg a)&=\neg\chi_U(a),&
 \chi_U(a\wedge b)&=\chi_U(a)\wedge\chi_U(b),\\
 \chi_U(a\vee b)&=\chi_U(a)\vee\chi_U(b).&&
\end{aligned}
\]
Here the operations on the right are in \(\mathbf2\). Conversely a
characteristic homomorphism has zero excluded, one included, meet closure
and upward closure (apply it to \(a\wedge b=a\)); its complement
equation gives decision, so it characterizes an ultrafilter.

For an ultrafilter \(U\), an element \(b\) lies in \(d(U)\) iff
\(\neg b\in U\), iff \(b\notin U\). Thus \(P=d(U)=B\setminus U\)
is a proper ideal. If \(a\wedge b\in P\), then
\(\neg a\vee\neg b\in U\), so one of \(\neg a,\neg b\) lies
in \(U\); hence one of \(a,b\) lies in \(P\).
Conversely a prime ideal \(P\) contains at least one of \(a,\neg a\),
because it contains their zero meet, and not both, because it excludes
their join \(1_B\). Its dual filter therefore decides complements and
is an ultrafilter; again it is \(B\setminus P\). Inclusion-preserving
duality takes maximal proper filters to maximal proper ideals. This also
proves that prime ideals are maximal, and that a maximal proper ideal is
prime by taking its dual filter. All families used are subsets of sets and
all witness combinations are finite; no choice function has been used.
\end{proof}
\end{annexitem}

\begin{annexitem}{annex:p01-bpi-quotient}{thm:p01-bpi-quotient}{Boolean quotient without representative choices}
\begin{proof}
First note the elementary identity
\[
 e\wedge(a\mathbin{\triangle}b)=0_B
 \quad\text{iff}\quad e\wedge a=e\wedge b.
\]
For the forward implication, distribute \(e\wedge a\) across
\(b\vee\neg b\) and \(e\wedge b\) across \(a\vee\neg a\).
The two difference terms vanish, leaving both sides equal to
\(e\wedge a\wedge b\). For the reverse implication substitute the
equal terms in each difference summand; it then contains the meet of an
element with its complement and is zero. In particular
\(\neg(a\mathbin{\triangle}b)\wedge a
=\neg(a\mathbin{\triangle}b)\wedge b\).

Reflexivity and symmetry of \(\sim_I\) follow from
\(a\mathbin{\triangle}a=0_B\) and symmetry of the displayed
expression. For transitivity, suppose
\(a\mathbin{\triangle}b,b\mathbin{\triangle}c\in I\) and set
\(e=\neg(a\mathbin{\triangle}b)\wedge
       \neg(b\mathbin{\triangle}c)\).
Then \(e\wedge a=e\wedge b=e\wedge c\). The identity implies
\(e\wedge(a\mathbin{\triangle}c)=0_B\), so
\[
 a\mathbin{\triangle}c\leq_B\neg e
  =(a\mathbin{\triangle}b)\vee(b\mathbin{\triangle}c)\in I.
\]
Downward closure gives transitivity.

If \(a\sim_I b\) and \(c\sim_I d\), use
\(e=\neg(a\mathbin{\triangle}b)\wedge
       \neg(c\mathbin{\triangle}d)\).
On meeting with \(e\), both inputs agree. Distribution shows that their
joins agree there, and associativity shows that their meets agree there.
The identity therefore bounds each of
\((a\vee c)\mathbin{\triangle}(b\vee d)\) and
\((a\wedge c)\mathbin{\triangle}(b\wedge d)\) by
\(\neg e=(a\mathbin{\triangle}b)\vee(c\mathbin{\triangle}d)\in I\).
Also \((\neg a)\mathbin{\triangle}(\neg b)=a\mathbin{\triangle}b\).
Thus \(\sim_I\) respects every operation.

Replacement forms the set of equivalence classes. Congruence makes
\([a]_I\wedge[b]_I=[a\wedge b]_I\), and similarly the join and
complement, single-valued independently of representatives. Each Boolean
axiom descends by applying the quotient map to that axiom in \(B\).
For any finitely many classes used in an axiom one may use representatives
in an existential argument; a global selector for all classes is neither
defined nor needed. The map \(q\) is a homomorphism and is surjective
by the definition of the quotient set. Since
\(a\mathbin{\triangle}0_B=a\), its zero kernel is exactly \(I\).
And \([0_B]_I=[1_B]_I\) iff
\(0_B\mathbin{\triangle}1_B=1_B\in I\); nontriviality is therefore
equivalent to properness of \(I\).

If \(Q\) is a prime ideal of the quotient, let \(P=q^{-1}[Q]\).
It contains zero, excludes one, is closed under joins, and is downward
closed because homomorphisms preserve the Boolean order. If
\(a\wedge b\in P\), then \(q(a)\wedge q(b)\in Q\);
primeness gives \(q(a)\in Q\) or \(q(b)\in Q\), as required.
For \(i\in I\), \(q(i)=[0_B]_I\in Q\), so \(I\subseteq P\).
\end{proof}
\end{annexitem}

\begin{annexitem}{annex:p01-bpi-finite-local}{thm:p01-bpi-finite-local}{Finite local Boolean evaluations}
\begin{proof}
Fix the finite list. Distributing
\(\bigwedge_{k=1}^{n}(b_k\vee\neg b_k)=1_B\) gives
\(\bigvee_{\varepsilon\in\{0,1\}^{n}}e_\varepsilon=1_B\).
If two indices differ in their \(k\)-th bit, the corresponding cells
have a meet containing \(b_k\wedge\neg b_k\), and hence have zero
meet. Discarding zero cells leaves a finite disjoint family with join one;
at least one survives since the algebra is nontrivial.

The joins of subfamilies of these cells form a subalgebra: joins take
unions of subfamilies, meets take intersections because the cells are
disjoint, and complement takes the remaining cells of the partition.
For the complement assertion, the selected and remaining joins have zero
meet and total join one, so uniqueness of complements gives the claim.
Each \(b_k\) is the join of cells whose \(k\)-th bit is one, by
distributing \(b_k\) against the other \(b_j\vee\neg b_j\).
Conversely every cell is a finite Boolean expression in the generators.
Thus the generated subalgebra is exactly this finite family of joins.
This is an algebraic distribution argument, not an application of Stone
representation or of prime-ideal existence.

Let \(E\subseteq I\) be finite and included among the generators.
The finite join \(j=\bigvee E\) lies in \(I\); since \(I\) is proper,
\(j\ne1_B\), so \(\neg j\ne0_B\). It belongs to the generated
subalgebra and is a join of cells. At least one cell in that join is
nonzero; fix such an \(e\). For any \(c\) in the subalgebra,
\(e\wedge c\) is either \(e\) or \(0_B\); hence precisely one of
\(e\leq_B c,e\leq_B\neg c\) holds. Meet and join of cell subfamilies
then show that \(h(c)=1\) iff \(e\leq_B c\) preserves both
operations, complements and constants. Every \(i\in E\) lies below
\(j\), while \(e\leq_B\neg j\), so \(e\wedge i=0_B\), and
\(h(i)=0\).

When \(n=0\), there is just the empty cell \(1_B\); the subalgebra
is \(\{0_B,1_B\}\). When \(E=\varnothing\), \(j=0_B\), and any
nonzero cell works. Fixing a witness for this one finite instance is an
existential inference in ZF; it is not simultaneous selection over the set
of all finite instances.
\end{proof}
\end{annexitem}

\begin{annexitem}{annex:p01-bpi-equivalence}{thm:p01-bpi-equivalence}{BPI, Boolean extension, and set ultrafilters}
\begin{proof}
\emph{Bare BPI implies ideal extension.} Given a proper ideal \(I\) of
\(B\), the quotient lemma gives a nontrivial algebra \(B/I\).
Apply BPI there and pull its prime ideal back along \(q\). The lemma
proves that the resulting proper prime ideal contains \(I\).
Conversely, \(\{0_B\}\) is a proper ideal of every nontrivial
Boolean algebra, so ideal extension implies bare BPI.

\emph{Ideal extension is equivalent to Boolean filter extension.}
The dual of a proper filter \(F\) is a proper ideal \(d(F)\).
If \(P\supseteq d(F)\) is prime, its dual
\(U=d(P)=B\setminus P\) is an ultrafilter and contains \(F\),
by inclusion-preserving duality. Conversely the dual of a proper ideal
extends to an ultrafilter \(U\); its dual prime ideal \(d(U)\)
contains the original ideal. The duality lemma proves every assertion
used in these constructions.

\emph{Boolean filter extension implies the set Ultrafilter Lemma.}
Apply it to \(\mathcal P(S)\). Its ultrafilter decides precisely one
of \(A,S\setminus A\) for every \(A\subseteq S\), by the duality
lemma. This is exactly the inherited formulation. If \(S\) is empty,
no proper filter was given and the instance is vacuous.

\emph{The set Ultrafilter Lemma implies ideal extension on every algebra.}
Let \(I\) be a proper ideal in a Boolean algebra \(B\), so
\(B\) is nontrivial. Let \(X=\mathbf2^B\) be the set of all
functions \(v:B\to\mathbf2\), not initially homomorphisms.
The constant-zero function shows that \(X\) is nonempty without Choice.
For \(b\in B\), put \(T_b=\{v\in X:v(b)=1\}\).
Use the following subsets of \(X\) as constraints:
\[
\begin{aligned}
 C_0&=\{v:v(0_B)=0\},& C_1&=\{v:v(1_B)=1\},\\
 C^{\neg}_b&=\{v:v(\neg b)=\neg v(b)\},&&\\
 C^{\wedge}_{a,b}&=\{v:v(a\wedge b)=v(a)\wedge v(b)\},&&\\
 C^{\vee}_{a,b}&=\{v:v(a\vee b)=v(a)\vee v(b)\},&&\\
 C_i&=\{v:v(i)=0\}&&(i\in I).
\end{aligned}
\]
Here every \(v\) is restricted to \(X\), and the operations on its
values are in \(\mathbf2\). Their collection is a set by Replacement
and Separation.

Every finite collection of constraints has nonempty intersection.
For each of its finitely many members fix one of the displayed constraint
indices defining it; such finite witness selection is proved in ZF by
induction on the number of members. Then collect all Boolean elements
occurring in those constraints,
including the operation results and any mentioned \(i\in I\), into
a finite list, and let \(E\) be its finite ideal constraints. The
finite local lemma gives a homomorphism on the generated subalgebra
annihilating \(E\). Define \(v\) to equal that homomorphism on the
subalgebra and to be zero outside it. This is a total function specified
by two disjoint cases and satisfies each of the selected constraints.
The empty collection is satisfied by any member of \(X\).
The generated set filter of all constraints is therefore proper.
Use the set Ultrafilter Lemma once to extend it to \(\mathcal U\).

By the duality lemma, its characteristic function
\(\chi:\mathcal P(X)\to\mathbf2\) is a Boolean homomorphism.
Define \(h(b)=\chi(T_b)\). If \(C\in\mathcal U\) and
\(C\cap A=C\cap A'\), then
\(\chi(A)=\chi(C)\wedge\chi(A)=\chi(C\cap A)
=\chi(C\cap A')=\chi(A')\).
Since every constraint belongs to \(\mathcal U\), their pointwise
identities imply
\[
\begin{aligned}
 h(0_B)&=0,& h(1_B)&=1,\\
 h(\neg b)&=\neg h(b),& h(a\wedge b)&=h(a)\wedge h(b),\\
 h(a\vee b)&=h(a)\vee h(b),& h(i)&=0\quad(i\in I).
\end{aligned}
\]
For example, on \(C^{\wedge}_{a,b}\) the sets
\(T_{a\wedge b}\) and \(T_a\cap T_b\) coincide; apply the preceding
equality and the homomorphism law for \(\chi\). On \(C^{\neg}_b\),
\(T_{\neg b}\) coincides with \(X\setminus T_b\). On \(C_0,C_1,C_i\),
the corresponding cylinder is respectively empty, all of the constraint,
or empty. This verifies each equation rather than presupposing a limiting
valuation.

The zero kernel \(P=\{b:h(b)=0\}\) is a proper ideal: homomorphism
laws give the constants and join closure, and preserve order for downward
closure. If \(h(a\wedge b)=0\), the two-element truth table forces
\(h(a)=0\) or \(h(b)=0\), so \(P\) is prime. It contains \(I\).
Only this one application of the set Ultrafilter Lemma goes beyond ZF.
In particular, no well-order of \(B\), extension of finite homomorphisms
to all of \(B\), family of preselected finite witnesses, or Boolean
representation theorem has been used.
\end{proof}
\end{annexitem}

\begin{annexitem}{annex:p01-bpi-propositional-compactness}{thm:p01-bpi-propositional-compactness}{The two compactness directions}
\begin{proof}
\emph{Ultrafilter extension implies compactness.} Let
\(\Gamma\) be finitely satisfiable over \(V\). The set
\(X=\mathbf2^V\) of total valuations is nonempty: if \(V\) is empty
its sole member is the empty function; otherwise the constant-zero
function is available. For a formula \(\varphi\), define its truth set
\(T_\varphi=\{v\in X:v(\varphi)=1\}\).
Structural recursion on finite formulas gives
\[
\begin{aligned}
 T_\bot&=\varnothing,& T_\top&=X,\\
 T_{\neg\varphi}&=X\setminus T_\varphi,&&\\
 T_{\varphi\land\psi}&=T_\varphi\cap T_\psi,&&\\
 T_{\varphi\lor\psi}&=T_\varphi\cup T_\psi.&&
\end{aligned}
\]
Finite satisfiability means precisely that every finite intersection of
truth sets for members of \(\Gamma\) is nonempty. Therefore they
generate a proper set filter. Extend it to \(\mathcal U\) by the
set Ultrafilter Lemma and let \(\chi\) be its characteristic
homomorphism on \(\mathcal P(X)\). Define a valuation
\(\nu(v)=\chi(T_{p_v})\) for \(v\in V\). Its values form one
function by Replacement, not separate choices. Induction on formula
construction using the displayed identities and the homomorphism laws
gives \(\nu(\varphi)=\chi(T_\varphi)\): the variable case is the
definition, both constant cases use the constant laws, and each of
negation, conjunction and disjunction uses the corresponding law.
Every \(T_\gamma\) lies in \(\mathcal U\), so
\(\nu(\gamma)=1\) for every \(\gamma\in\Gamma\). This supplies
compactness. The equivalence theorem supplies the same implication
from BPI without importing AC.

\emph{Compactness implies Boolean ideal extension.} Let \(I\) be a
proper ideal of a nontrivial Boolean algebra \(B\). Introduce distinct
variables \(p_b\) for \(b\in B\), using tagged syntax if an element
of \(B\) already happens to be a formula. Take the following set of
finite formulas:
\[
\begin{gathered}
 \neg p_{0_B},\qquad p_{1_B},\qquad
 p_{\neg b}\Leftrightarrow\neg p_b,\\
 p_{a\wedge b}\Leftrightarrow(p_a\land p_b),\qquad
 p_{a\vee b}\Leftrightarrow(p_a\lor p_b),\qquad
 \neg p_i\quad(i\in I).
\end{gathered}
\]
Here \(\Leftrightarrow\) abbreviates the Boolean expression for equal
truth values; it is formula equivalence, not elementary equivalence of
structures. For any finite subset, collect all mentioned algebra
elements, including operation results, into a finite list. Apply the
finite local lemma to its ideal constraints and define a total assignment
to be its finite-subalgebra homomorphism there and zero on all other
elements. Every selected formula is then true. Thus the displayed set is
finitely satisfiable, and compactness gives a valuation satisfying all
of it. The function \(h(b)=\nu(p_b)\) consequently preserves every
Boolean operation and both constants, and annihilates \(I\). Its
zero kernel is a prime ideal containing \(I\), by the kernel argument
in the preceding annex item. The equivalence theorem completes all the
other implications.

Finally, this proof uses compactness for the variable set indexed by all
of \(B\), not just a countable subset. For a fixed enumerable alphabet,
the separate binary-tree argument of Part~6 uses its supplied enumeration
and the fixed order of the two bits and requires no ultrafilter principle.
The two assertions have different domains. Truth constants cause no
extra strength: for a nonempty variable set fix one variable \(p\)
and replace \(\top,\bot\) by \(p\lor\neg p,p\land\neg p\).
This translation preserves truth under every valuation. For an empty
alphabet the single empty valuation evaluates every constants-only
formula; finite satisfiability already excludes any false formula by
testing its singleton. Thus the constants convention does not hide a
choice assumption or a failed empty-alphabet case.
\end{proof}
\end{annexitem}

\begin{annexitem}{annex:p01-bpi-ac-implication}{thm:p01-bpi-ac-implication}{The precise Zorn application}
\begin{proof}
Assume AC and fix a proper Boolean filter \(F\subseteq B\).
The set
\(\mathcal Q=\{G\subseteq B:G\text{ is a proper filter and }F\subseteq G\}\)
exists by Separation from \(\mathcal P(B)\) and is nonempty because it
contains \(F\). Order it by inclusion. A nonempty chain
\(\mathcal C\) has upper bound \(G_*=\bigcup\mathcal C\): it
contains \(F\) and \(1_B\), excludes \(0_B\), and is upward closed.
If \(a,b\in G_*\), they belong to two members of the chain, and
one of those contains the other. Both elements and their meet therefore
belong to that larger filter and hence to \(G_*\). It is a proper
filter. The empty chain has the upper bound \(F\), so the hypothesis
including empty chains is satisfied. The Choice-to-Zorn theorem of Part~6
now gives a maximal \(U\in\mathcal Q\). Any proper filter extending
\(U\) still contains \(F\), so maximality in \(\mathcal Q\)
is maximality among all proper Boolean filters. Thus \(U\) is an
ultrafilter. The ZF equivalence already proved gives BPI.
The only non-ZF premise in this argument is Zorn supplied by AC; the
chain union and all earlier algebraic constructions use ZF alone.
\end{proof}
\end{annexitem}

\begin{annexitem}{annex:p01-bpi-set-boundaries}{thm:p01-bpi-set-boundaries}{Finite and cofinite filters}
\begin{proof}
For any set \(S\), if an ultrafilter is generated by a single
\(A\subseteq S\), properness gives \(A\ne\varnothing\). Fix
\(s\in A\). If \(A\ne\{s\}\), then neither \(\{s\}\) nor
\(S\setminus\{s\}\) contains \(A\), so neither is in the
generated filter, contradicting ultrafilter decision. Thus \(A=\{s\}\),
and the filter is \(U_s\). Conversely every point filter is generated
by its singleton and is an ultrafilter by decision. This verifies the
equivalence of the two principal-ultrafilter conventions.

Suppose \(S\) is finite and \(F\) is a proper filter on it.
Then \(S\ne\varnothing\) and \(\mathcal P(S)\), hence \(F\),
is finite in ZF. These finite-set facts follow by induction on a finite
enumeration of \(S\): adjoining one point doubles the power-set list,
and a subset of a finite list has a finite list. The intersection
\(A=\bigcap F\) is therefore a finite intersection of members of
\(F\), belongs to \(F\), and is nonempty by properness. Fix one
\(s\in A\). Then \(U_s=\{Y\subseteq S:s\in Y\}\) is a proper
filter deciding every subset by whether it contains \(s\), hence an
ultrafilter, and it contains \(F\). Selecting this single witness is
not an axiom of simultaneous choice for a family.

If \(U\) is any ultrafilter on finite \(S\), the same finite
intersection argument gives a point \(s\in\bigcap U\) and
\(U\subseteq U_s\). Maximality gives equality. More directly,
primeness applied repeatedly to the finite union of singletons making
\(S\) gives a singleton in \(U\); two distinct singletons cannot
both belong to a proper filter. Thus the principal point is unique.

Now let \(S\) be infinite, meaning not finite, without assuming a
countably infinite subset. Its cofinite family
\(F_{\mathrm{cof}}=\{Y\subseteq S:S\setminus Y\text{ is finite}\}\)
contains \(S\) and excludes \(\varnothing\). It is upward closed
because a subset of a finite set is finite. It is meet closed because
the union of two finite complements is finite, a fact proved by
concatenating finite lists and deleting repetitions. Therefore it is
a proper filter in ZF. Assuming BPI, the equivalence theorem extends it
to an ultrafilter \(U\). If \(T\subseteq S\) is finite, then
\(S\setminus T\in F_{\mathrm{cof}}\subseteq U\), so
\(T\notin U\), since the two have empty intersection. In particular
no singleton belongs to \(U\), while \(U_s\) contains \(\{s\}\)
for every \(s\in S\). Thus \(U\) is not principal. The extra
existence step is exactly ultrafilter extension, not a purported ZF
construction of a free ultrafilter.
\end{proof}
\end{annexitem}
